\documentclass[10pt]{article}
\usepackage[margin=0.85in]{geometry}
\usepackage[T1]{fontenc}
\usepackage{lmodern}
\usepackage{microtype}
\usepackage{enumitem}
\setlist{nosep,leftmargin=1.2em}
\usepackage[hidelinks]{hyperref}
\setlength{\parskip}{0.35em}
\setlength{\parindent}{0pt}
\pagestyle{plain}

\begin{document}
{\large\bfseries Advisor brief: the revised combined empirical paper and a proposed method scope}\\[0.2em]
{\small 16 October 2026. Revision of the combined manuscript at commit \texttt{3b5065805}; the later
\texttt{PCRL\_Combined\_Empirical\_Revised} files were not available on this machine and were not revised.}

\section*{1. What the paper now asks}
For stored representation models and published defenses, what can a recipient recover about a declared protected
attribute from everything it receives, alone or linked with one other recipient, and how useful is that release for
the recipient's task? The paper now argues from the problem to the evidence. Protection methods make promises scoped to
an object and an attacker class: LEACE covers linear classifiers reading the representation, and FARE certifies
classifiers that read its released cell. A recipient may receive more than that object, namely a task output and, if
recipients link records, another recipient's output. The results separate three outcomes: a method failing its own
test, a test holding while recovery occurs outside its scope, and recovery through a surface the test does not cover.

\section*{2. Why this adversary and this access}
The claims concern a release policy, so the adversary holds what the policy delivers. An authorized recipient
necessarily possesses its own release, so that part needs no further justification. Two capabilities are granted, not
established. \emph{Linkage}: two recipients serving the same person may pool releases by record; we grant perfect
linkage, the attacker-favourable limit, without assuming that governance prevents linkage or that pooling is
unavoidable. \emph{Labelled auxiliary data}: our attackers are supervised and learn from releases of other records with
the attribute, which models a recipient that knows the attribute for some of the people it serves; availability, size
and population match are assumptions. A stronger attacker alone does not justify more access. Inputs held by the
recipient, published parameters, repeated queries and insiders are distinct threat models, not necessarily breaches;
they are untested, not dismissed. The paper now has a table that states, for each experiment, what the recipient
receives, the labelled sample, how attackers are fitted, selected and scored, what they know about the defense, how
releases are linked and what is untested.

\section*{3. What the results establish}
\begin{itemize}
\item \textbf{Scope, not failure.} Official LEACE meets its linear bound on every map, while nonlinear attackers recover
the attribute about as well as from untreated features (AUC 0.817 and 0.863 against 0.815 and 0.870).
\item \textbf{A method's own test failing.} The encoders' stored linear check fails on 12 of 42 attribute--seed
combinations.
\item \textbf{Other surfaces.} Releasing the untreated model's output beside defended features restores recovery to
0.76--0.79, and outputs reveal more than the true label alone would.
\item \textbf{Utility changes the reading.} With useful heads, FARE's score recovers 0.539 against 0.774 untreated, but a
same-budget tree without the fairness term already reaches 0.648, and FARE's accuracy is not shown to stay within one
point. On an easier task, compression reproduces most of FARE's effect.
\item \textbf{Linked recipients.} Two linked recipients recover race better than either alone in every output format,
including hard decisions (0.611 against 0.576).
\end{itemize}
All of this is development evidence on previously used rows, with intervals conditional on the fitted heads and
attackers.

\section*{4. What remains unproved}
There is no population privacy guarantee: recovery is a property of our attacker slate and data budget. Differences in
recovery may reflect information removed, limits of our attackers, or both, and the paper does not separate them. A
deterministic release need not be reversible, and the paper does not establish whether any release keeps or discards
the attribute's posterior. The excluded threat models, imperfect linkage and other auxiliary budgets are untested.
Nothing is fresh confirmation, and there is no new method and no composition or synergy result.

\section*{5. Proposed objective for the future method}
Develop purpose-specific releases that retain the usefulness required by each permitted task while reducing
sensitive-attribute recovery from everything recipients receive, both individually and when releases are linked. The
endpoint is empirical: recovery under declared access, attacker families and auxiliary-data budgets, reported beside the
utility of what is released. A mechanism may aim to remove information that the tasks do not need and still be judged
by finite attacks; the two goals are compatible. A claim about arbitrary attackers would need a theorem; exact examples
demonstrate their own laws, not real-population privacy. A competitive method need not be a joint design, and joint
benefit, competitiveness and novelty are separate questions. Existing utility caps, release contracts and references
stay unchanged. No competitive new mechanism has been established, and agreeing on this scope does not supply one.

\section*{6. Feedback and decisions needed}
\begin{enumerate}
\item Is the question and the stated threat model the right framing for the empirical paper, and is the per-experiment
access table sufficient?
\item Should the paper stay at its current evidence cutoff? Possible extensions are listed separately and none is
proposed for immediate work.
\item What usefulness does the service actually need: calibrated probabilities, ranking scores, or decisions? This
choice defines the method problem.
\item Which attacker families, auxiliary budgets, linkage model and comparators should a future method be judged
against, and is a formal guarantee a goal?
\end{enumerate}
\end{document}
